\chapter{Brainteaser：把陌生题压缩成可证明结构}
\label{ch:lower-brainteasers}

\begin{itemize}
  \input{tex/generated/prerequisites/lower-brainteasers}
  \item 学习产出：能把陌生题压缩成状态、不变量、下界和构造，并在有限时间内给出可核验的口述或笔试答案。
  \item 研究边界：谜题表现不能替代方向模型、样本外验证与交易摩擦证据。
\end{itemize}

Brainteaser 不是猜谜比赛。量化研究笔面试真正考察的是：能否在信息不完整、时间有限时，迅速选择状态、写出不变量、给出下界，并说明答案依赖哪些假设。本章把这类能力拆成五种可复用动作；题目来自《量化金融面试实用指南》第二章的核心能力，但题面、参数和解答均重新设计。

\section{先缩小状态，再开始计算}

遇到海盗分配、动物过河或多人博弈，第一问不是“有哪些招”，而是“决定未来的最小状态是什么”。若历史只通过当前人数、轮次或余数影响以后，就把它压缩成状态 $s$；若目标是保证成功，则从最小规模向后归纳。

\begin{example}
五位研究员按资历提议分配奖金，方案通过需要至少一半票数，否决者离场。不要枚举所有分配；从一人、两人局面逆推，每一步先识别哪些人的下一轮收益最低，再用最少成本买票。
\end{example}

Python 中可以递归枚举小规模局面，但面试答案必须先说明状态、终止条件和参与者的偏好顺序；否则程序只是在替一个未定义的博弈求解。

\section{逻辑约束：把语言变成可排除的集合}

真假标签、称球和过桥题适合先声明概率空间 $(\Omega,\mathcal F,\mathbb P)$，再把可行状态写成候选集合 $\Omega$。一次观察把它切成若干分支；最坏分支的大小决定信息效率。若一次实验最多有 $b$ 种结果，$k$ 次实验区分 $N$ 个状态至少需要 $b^k\ge N$。这个下界先告诉你“几次不可能”，再指导构造。

\begin{example}
若天平一次有左重、平衡、右重三种结果，两次称量最多编码九个叶节点。因此想在十二种互斥状态中保证定位，两次称量从信息量上已不可能。
\end{example}

\section{不变量、对称与奇偶性}

硬币堆、变色球、开关和蚂蚁题经常故意诱导逐步模拟。更短的路线是寻找每次操作都保持的量：总和模 $m$、逆序数奇偶性、颜色计数差，或在交换标签后不变的轨迹集合。对称不是“看起来一样”，而是存在保持概率或可达性的双射。

\begin{example}
若每次操作同时翻转两个开关，则“亮灯数的奇偶性”保持不变。于是从偶数盏亮灯出发，不可能到达奇数盏亮灯的目标；这比搜索所有开关序列更强。
\end{example}

\section{条件概率与期望：先声明信息集}

生日、卡牌与错标袋问题的陷阱通常不是算术，而是条件事件含糊。先写样本空间和已知信息，再计算 $P(A\mid B)$；若过程有多轮，优先尝试全期望公式或指示变量线性性。不要因为随机变量相关，就误以为其指示变量不能逐项求期望。

\section{递推、博弈与停止规则}

轮流取物、囚徒策略和序列题要区分三件事：状态转移、双方信息、停止条件。先计算小规模表格寻找模式，再用归纳或模运算证明；小规模程序只是猜想生成器，不是证明。

\section{精选题：五种方法}

以下每题只有一个主归属，交叉章节只可引用，不重复计入覆盖率。四级训练覆盖\textbf{口述概念、笔试推导、数值编程、研究判断}；建议先口述建模，再在纸上给出最坏情形证明。

\subsection*{口述概念}
限时解释状态、信息与关键不变量：GB-2-01、02、05、11--14、17--19、21、23、26、33、35、36。

\subsection*{笔试推导}
给出下界、构造与完备性证明：GB-2-03、04、06--08、15、16、20、24、25、27、28、30--32、34。

\subsection*{数值编程}
将闭式解或统计量恢复写成可验证程序：GB-2-10、22。

\subsection*{研究判断}
在信息、成本或协议限制下评估结论是否成立：GB-2-09、29。

下面按逆向归纳、信息下界、不变量、条件概率和递推五种方法组织精选题。来源映射和完整覆盖台账保存在附录与项目报告中；读者首先需要看到的是题目要求掌握的思路。

\begin{enumerate}[leftmargin=*]
  \item \MFQInterviewMeta{口述}{12 分钟}{难}{逆向归纳与博弈均衡}{GB-2-01 Screwy pirates}六位合伙人依次提出整数奖金分配；提案需严格过半，否决者退出。所有人依次偏好生存、奖金、让别人退出。写出从一人局面开始的逆推算法，并说明平票规则改变时哪里必须重算。
  \item \MFQInterviewMeta{口述}{8 分钟}{中}{状态机与安全策略}{GB-2-02 Tiger and sheep}一只守卫和两只闯入者在三节点通道上轮流移动。把题抽象成有限状态图，说明怎样判定守卫是否存在必胜策略，而不是逐句描述追逐过程。
  \item \MFQInterviewMeta{笔试}{12 分钟}{中}{约束搜索与下界}{GB-2-03 River crossing}一名研究员要把模型文件、密钥和原始数据送过河。船每次只能载研究员和一件物品，且必须由研究员驾驶；研究员不在某岸时，密钥不得与原始数据同岸，模型文件也不得与密钥同岸。给出最短运输序列，并用状态图的层数证明不能再少一次渡河。
  \item \MFQInterviewMeta{笔试}{10 分钟}{中}{同余与日历建模}{GB-2-04 Birthday problem}某人的生日在连续四年中分别落在周二、周三、周四、周六。只用闰年规则判断哪一年跨过了闰日，并列出仍无法确定的信息。
  \item \MFQInterviewMeta{口述}{10 分钟}{中}{不完全信息与策略}{GB-2-05 Card game}两人各抽一张 $1$ 到 $9$ 的牌，只看自己的牌，可选择加注或弃牌。说明为什么“按牌面单调加注”是合理候选策略，并指出还需哪些收益参数才能求阈值。
  \item \MFQInterviewMeta{笔试}{8 分钟}{中}{构造与误差边界}{GB-2-06 Burning ropes}两根绳子的总燃烧时间各为一小时，但燃速不均匀。设计测量四十五分钟的方法，并说明方案依赖的唯一物理假设。
  \item \MFQInterviewMeta{笔试}{15 分钟}{难}{信息论下界与决策树}{GB-2-07 Defective ball}十二个外观相同的零件中恰有一个重量异常，且未知偏轻偏重。先用叶节点计数给出称量次数下界，再说明构造决策树时必须保留哪些状态。
  \item \MFQInterviewMeta{笔试}{10 分钟}{中}{估值与计数}{GB-2-08 Trailing zeros}不用计算阶乘本身，求 $125!$ 末尾零的个数，并推广到任意进制 $b$；指出只数 $5$ 的因子何时会失效。
  \item \MFQInterviewMeta{研究判断}{12 分钟}{难}{分组淘汰与比较下界}{GB-2-09 Horse race}二十五个策略只能五个一组回测，每次只得到组内相对名次。没有计时器时找出全局前三至少需要几轮？给出候选淘汰表。
  \item \MFQInterviewMeta{数值编程}{12 分钟}{难}{递推与闭式验证}{GB-2-10 Infinite sequence}序列满足 $a_{n+2}=3a_{n+1}-2a_n$，只给定 $a_0,a_1$。推导闭式、验证极限条件，并说明数值递推何时会放大舍入误差。
  \item \MFQInterviewMeta{口述}{10 分钟}{中}{几何装箱与反例}{GB-2-11 Box packing}判断边长为 $1$ 的立方体能否放入尺寸为 $0.9\times1.2\times1.4$ 的长方体。该长方体体积大于 $1$、空间对角线也大于 $\sqrt3$；证明它仍无法容纳单位立方体，并解释为什么体积和对角线都不是充分条件。
  \item \MFQInterviewMeta{口述}{10 分钟}{中}{位值编码与可辨识性}{GB-2-12 Calendar cubes}用两个立方体的六个面标数字，使每天都能并排显示 $01$ 到 $31$。先写出数字出现次数的必要条件，再给出一组可行标记，并说明为什么 $6$ 与 $9$ 可共用一个旋转面。
  \item \MFQInterviewMeta{口述}{8 分钟}{中}{最坏情形决策}{GB-2-13 Door to offer}三扇门后分别是录用、加面和拒绝。证明一次是/否询问不能保证识别三种结果；若允许第二个根据首答自适应选择的是/否问题，给出最少询问协议和完整决策树。
  \item \MFQInterviewMeta{口述}{10 分钟}{中}{通信协议与公共知识}{GB-2-14 Message delivery}两名研究员只能通过一名可能丢失消息但不会篡改消息的信使通信。解释为什么“我知道你收到”会产生无限确认链，并给出工程上可终止的超时、序号与重试协议及其保证边界。
  \item \MFQInterviewMeta{笔试}{10 分钟}{中}{递推与奇偶性}{GB-2-15 Last ball}袋中有红蓝球，每次取出两球：若同色则放回一个蓝球，若异色则放回一个红球。设初始蓝球数为 $B$ 且总球数为 $N$，用不变量判断最后一球的颜色；不要模拟抽取顺序。
  \item \MFQInterviewMeta{笔试}{10 分钟}{中}{奇偶不变量}{GB-2-16 Light switches}一排 $n$ 个开关初始全灭，每次只能翻转相邻两个。刻画所有可达状态，并给出到达某个目标状态的线性时间判定法。
  \item \MFQInterviewMeta{口述}{8 分钟}{中}{期望值与条件缺失}{GB-2-17 Quant salary}公司给出“首年奖金期望为固定工资的 $40\%$”。说明仅凭期望为什么不能比较两个 offer 的效用，并列出至少四项还需询问的分布、约束或相关性信息。
  \item \MFQInterviewMeta{口述}{8 分钟}{中}{对称变换}{GB-2-18 Coin piles}一堆硬币中已知恰有 $k$ 枚正面，但不能观察单枚朝向。蒙眼把它分成两堆，使两堆正面数相同；证明翻转操作为何成立。
  \item \MFQInterviewMeta{口述}{8 分钟}{易}{条件排除}{GB-2-19 Mislabeled bags}三个袋子标签全部贴错，分别装两种资产及其混合。最少抽取几次可恢复全部标签？写出第一次抽样后为何没有分支歧义。
  \item \MFQInterviewMeta{笔试}{12 分钟}{难}{公共知识与递推}{GB-2-20 Wise men}三人各戴一顶黑或白帽，只能看见别人。主持人宣布“至少一顶黑帽”。若前两人依次说不知道，第三人如何推断？明确公共信息如何逐轮更新。
  \item \MFQInterviewMeta{口述}{8 分钟}{中}{圆周切割与最坏间隔}{GB-2-21 Clock pieces}从钟面上选两条经过圆心的直线切成四块。能否让四块上的数字和相等？先用总和给出目标，再用对径数字配对证明或否定可行性。
  \item \MFQInterviewMeta{数值编程}{10 分钟}{中}{等差和与缺失量}{GB-2-22 Missing integers}从 $1,\ldots,n$ 中删去两个不同整数，只给剩余数的和与平方和。推导恢复两数的公式，并写出输入不一致的判定条件。
  \item \MFQInterviewMeta{口述}{10 分钟}{难}{编码与称量}{GB-2-23 Counterfeit coins I}若一批硬币中假币都偏轻且数量未知，只有一次精密秤读数，怎样设计不同取样权重来编码假币集合？讨论量程和误差限制。
  \item \MFQInterviewMeta{笔试}{6 分钟}{易}{鸽巢原理}{GB-2-24 Matching socks}黑、灰、蓝三色袜子混放，闭眼最少取几只才能保证一双同色？若要求保证两双且两双颜色不同，答案如何变化？
  \item \MFQInterviewMeta{笔试}{8 分钟}{中}{图的度数和}{GB-2-25 Handshakes}一场活动中每两人至多握手一次。证明握手次数等于度数和的一半，并推出奇数度人数必为偶数。
  \item \MFQInterviewMeta{口述}{8 分钟}{中}{鸽巢原理与补图}{GB-2-26 Have we met before?}六位候选人中证明必有三人两两认识，或三人两两不认识。先固定一人，再对其五条关系分类。
  \item \MFQInterviewMeta{笔试}{8 分钟}{中}{轨迹对称}{GB-2-27 Ants on a square}单位正方形的四条边是轨道，四个顶点都是出口。若干蚂蚁从边上给定位置出发，以相同速度 $v$ 按给定顺/逆时针方向移动，相撞即掉头。用每只蚂蚁沿初始方向到下一个顶点的距离表示全部离开的时间，并证明“相撞掉头”等价于“穿透而过”的标签交换。
  \item \MFQInterviewMeta{笔试}{15 分钟}{难}{三进制编码}{GB-2-28 Counterfeit coins II}一次天平称量有三种结果。给定三次称量，设计可识别的异常状态编码，并说明为何每枚物品的编码不能与另一枚互为相反数。
  \item \MFQInterviewMeta{研究判断}{15 分钟}{难}{模运算与协同策略}{GB-2-29 Prisoner problem}多人依次观察他人的二元状态 $0/1$，每人回答时只能公开一个比特。设计奇偶校验协议，让除至多一人外都能恢复自己的状态；说明牺牲者的公开比特保存了什么信息。
  \item \MFQInterviewMeta{笔试}{8 分钟}{中}{十进制与模 $9$}{GB-2-30 Division by 9}证明一个十进制整数与其各位数字和模 $9$ 同余；设计可检测单个数字抄错的校验，并说明哪些错误模式仍无法被发现。
  \item \MFQInterviewMeta{笔试}{12 分钟}{难}{模不变量与可达性}{GB-2-31 Chameleon colors}初始有红、绿、蓝变色龙各 $8,4,3$ 只。每次两只不同颜色的变色龙相遇，二者都变成第三种颜色。用两两颜色计数之差模 $3$ 的不变性，判定是否可能最终全部同色。
  \item \MFQInterviewMeta{笔试}{10 分钟}{中}{构造与归纳}{GB-2-32 Coin split problem}有 $2n$ 枚硬币，其中 $n$ 枚正面朝上但位置未知。蒙眼把它们分成两堆，使两堆正面数相同；给出构造并证明它与初始排列无关。
  \item \MFQInterviewMeta{口述}{8 分钟}{易}{操作下界}{GB-2-33 Chocolate bar problem}一块 $m\times n$ 的巧克力每次只能拿一块现有碎片掰成两块。分成 $mn$ 个小格最少需要几次？证明下界并给出达到下界的任意策略。
  \item \MFQInterviewMeta{笔试}{12 分钟}{中}{相对速度与周期}{GB-2-34 Race track}两名跑者在长度 $L$ 的环形跑道同点同向出发，速度分别为 $u>v>0$。推导第 $k$ 次追及的时刻与位置；再说明反向出发时公式如何改变。
  \item \MFQInterviewMeta{口述}{10 分钟}{中}{反证法与构造}{GB-2-35 Irrational number}证明存在两个无理数 $a,b$ 使 $a^b$ 为有理数。答案不得预先假定 $\sqrt2^{\sqrt2}$ 的有理性，而应按其是否有理分情况构造。
  \item \MFQInterviewMeta{口述}{12 分钟}{难}{编码、信息与最坏保证}{GB-2-36 Rainbow hats}一列研究员各戴有限种颜色之一的帽子，只能看见前方并依次回答自己的颜色。设计使最坏情况下错误人数有界的编码协议，并明确答案依赖颜色数、回答顺序及能否听见先前回答。
\end{enumerate}

\section{作答协议}

每道题的完整答案都应包含：一句建模、一个下界或不变量、构造步骤、边界假设和至少一个常见误区。高频题还要准备面试官追问：参数变化后哪些结论仍成立，哪些必须重算。

本章不设独立综合项目；它的学习目标是把“状态--不变量--下界--构造--限制”这套方法带入六条研究路线，并为各路线的 Capstone 提供可迁移的推理习惯。谜题表现不能替代方向模型、样本外检验与交易摩擦证据。
